\subsection{\'Etale descent and reduced norms}\label{sec:paper-local-norm}
The main result of this subsection is Theorem~\ref{thm:paper-local-norm}, which identifies the positive-degree $K$-groups of a central division algebra over a nonarchimedean local field with those of the field, compatibly with splitting transfers and field norms.
We first fix the conventions on norms and algebras that are used in the rest of the paper.
Let $R\to S$ be a homomorphism of rings such that $S$ is perfect as a right $R$-module.
Restriction of scalars $\operatorname{Perf}(S)\to\operatorname{Perf}(R)$ induces the \emph{norm}
\[
 N_{S/R}\colon K(S)\longrightarrow K(R).
\]
The norms commute with flat base change: if $R\to S$ is a homomorphism of algebras over a commutative ring $k$ and $k'$ is a flat commutative $k$-algebra, then $N_{S\otimes_kk'/R\otimes_kk'}\circ(-\otimes_kk')=(-\otimes_kk')\circ N_{S/R}$.
For composable maps $R\to S\to T$ with $S$ perfect as a right $R$-module and $T$ perfect as a right $S$-module, we have $N_{T/R}=N_{S/R}\circ N_{T/S}$.

Let $F$ be a field, let $A$ be a central simple algebra over $F$, and let $D$ be the division algebra in the Brauer class of $A$.
Since Morita equivalence identifies $K(A)$ with $K(D)$, we work with $D$ instead of $A$ from now on.
We write $i_D\colon K(F)\to K(D)$ for extension of scalars.
We write $\Nrd_i\colon K_i(D)\to K_i(F)$ for the reduced norms ($i=1,2$).
We write $F^{\mathrm{sep}}$ for a separable closure of $F$.
For a field extension $E/F$, put $D_E=D\otimes_FE$, a central simple $E$-algebra that need not be a division algebra, and let $D'$ be the division algebra in its Brauer class.
When we apply a result stated for central division algebras to $D_E$, we identify $K(D_E)$ with $K(D')$ by the Morita equivalence given by a simple $(D',D_E)$-bimodule with compatible $E$-actions.
Norms are still taken for $D_E$ itself: if $D_E\simeq M_r(D')$, then $N_{D_E/E}=r\,N_{D'/E}$ under this identification.
If $E/F$ is finite and splits $D$, we identify $K(E)$ with $K(D_E)$ by the Morita equivalence given by a simple $(E,D_E)$-bimodule on which the two actions of $E$ agree; this identification does not depend on the choice of the simple bimodule.
The splitting transfer is the composite
\[
 \theta_E\colon K(E)\simeq K(D_E)\xrightarrow{\ N_{D_E/D}\ }K(D).
\]
A simple $(E,D_E)$-bimodule has dimension $\deg D$ over $E$, so $N_{D_E/E}=\deg D$ under this identification, and transitivity of norms gives
\[
 N_{D/F}\circ\theta_E=\deg D\cdot N_{E/F}.
\]

Throughout this subsection, $F$ is a nonarchimedean local field with residue field $k$ of characteristic $p$.

\begin{theorem}\label{thm:paper-local-norm}
Let $F$ be a nonarchimedean local field, and let $D$ be a central division algebra over $F$.
For each $j\ge1$ there is an isomorphism
\[
 \nu_{j,D}\colon K_j(D)\xrightarrow{\ \sim\ }K_j(F)
\]
satisfying $\nu_{j,D}\theta_E=N_{E/F}$ for every finite splitting field $E/F$ of $D$.
For a finite separable extension $F'/F$ it satisfies $\nu_{j,D_{F'}}\circ\res_{F'/F}=\res_{F'/F}\circ\nu_{j,D}$ and $\nu_{j,D}\circ N_{D_{F'}/D}=N_{F'/F}\circ\nu_{j,D_{F'}}$, with the convention above for $D_{F'}$, and $\deg D\cdot\nu_{j,D}=\pi_jN_{D/F}$.
In degrees $1$ and $2$ it is $\Nrd_1$ and $\Nrd_2$, respectively.
\end{theorem}
The proof of Theorem~\ref{thm:paper-local-norm} is given at the end of this subsection.

Let $K^{\et}(D)$ be the spectrum of global sections of the hypersheafification of $L\mapsto K(D_L)$ on the small \'etale site of $F$, and let
\[
 c_D\colon K(D)\longrightarrow K^{\et}(D)
\]
be the \'etale comparison map.
For a finite separable extension $E/F$, the commutative \'etale comparison $c_E\colon K(E)\to K^{\et}(E)$ induces isomorphisms on $\pi_i$ for $i\ge0$ \cite[Theorems 7.12(2) and 7.13]{ClausenMathew}.
At each prime $\ell$ the dimension bound in that theorem is at most two: it is $\vcd_\ell(E)$ for $\ell\ne\operatorname{char}E$, and it is $1+\log_p[E:E^p]=2$ for $\ell=p=\operatorname{char}E$.
Write $G_F=\Gal(F^{\mathrm{sep}}/F)$ and $G_k=\Gal(k^{\mathrm{sep}}/k)$.
The proof of Theorem~\ref{thm:paper-local-norm} compares $K(D)$ with its \'etale sheafification, whose geometric stalks agree with those of $F$ by Morita equivalence.
Lemma~\ref{lem:paper-local-complete-high} shows that the comparison is an isomorphism after completion in degrees at least two.
Lemma~\ref{lem:paper-local-galois-vanishing} gives the vanishing range of the Galois cohomology that enters integral descent, and Lemma~\ref{lem:paper-local-low-descent} uses it to show that $\pi_0$ of the \'etale theory is torsion-free and that the comparison on $\pi_1$ is $\Nrd_1$.
Lemma~\ref{lem:paper-local-compare} combines Lemmas~\ref{lem:paper-local-complete-high} and~\ref{lem:paper-local-low-descent} into an integral comparison, and Lemma~\ref{lem:paper-local-even-coinvariants} describes the even degrees by coinvariants.

\begin{lemma}\label{lem:paper-local-complete-high}
For every prime $\ell$, the map
\[
 (c_D)^\wedge_\ell\colon K(D)^\wedge_\ell\longrightarrow K^{\et}(D)^\wedge_\ell
\]
induces an isomorphism on $\pi_j$ for every $j\ge2$.
If $\ell\ne\operatorname{char}F$, the localization unit
\[
 K(D)^\wedge_\ell\longrightarrow L_{K(1),\ell}K(D)
\]
induces an isomorphism on $\pi_j$ for every $j\ge2$.
\end{lemma}
\begin{proof}
We first compute the targets after derived $\ell$-completion.
Let $\ell\ne\operatorname{char}F$ and $a\ge1$.
By Morita equivalence, the geometric stalks are
\[
 \pi_t\bigl(K(D_{F^{\mathrm{sep}}})/\ell^a\bigr)\simeq\pi_t\bigl(K(F^{\mathrm{sep}})/\ell^a\bigr)\simeq
 \begin{cases}
  \Z/\ell^a(b), & t=2b\ge0,\\
  0, & \text{otherwise},
 \end{cases}
\]
by Gabber--Suslin rigidity and the computation of $K(F^{\mathrm{sep}})/\ell^a$ \cite{SuslinClosed}, and for $L_{K(1),\ell}K(D_{F^{\mathrm{sep}}})/\ell^a$ the same formula holds with $\Z/\ell^a(b)$ in every even degree $t=2b\in\Z$.
The height-one localization satisfies \'etale hyperdescent, since it is the value of a weakly localizing invariant with tensor coefficients $\operatorname{Perf}(D)$ \cite[Example 5.12 and Theorem 7.14]{ClausenMathew}.
The descent spectral sequences therefore have
\[
 E_2^{s,2b}=H^s(F,\Z/\ell^a(b)),\qquad E_2^{s,t}=0\quad(t\text{ odd}).
\]
Since the geometric degrees are even and $\cd_\ell(F)\le2$, the differentials vanish, and since these groups are finite, their inverse systems are Mittag--Leffler.
By hyperdescent of the height-one localization, the localization unit factors as $K(D)^\wedge_\ell\xrightarrow{(c_D)^\wedge_\ell}K^{\et}(D)^\wedge_\ell\to L_{K(1),\ell}K(D)$.
The descent calculations for $K^{\et}(D)/\ell^a$ and $L_{K(1),\ell}K(D)/\ell^a$ agree in positive degrees, so, after passage to the inverse limit, the second map induces isomorphisms on $\pi_j$ for $j\ge1$.
Thus the two assertions of the lemma are equivalent for $\ell\ne\operatorname{char}F$.
Passing to the inverse limit over $a$ gives, for $m\ge1$,
\[
 \pi_{2m-1}\bigl(K^{\et}(D)^\wedge_\ell\bigr)\simeq H^1(F,\Z_\ell(m)),\qquad
 \pi_{2m}\bigl(K^{\et}(D)^\wedge_\ell\bigr)\simeq H^2(F,\Z_\ell(m+1));
\]
the other term in degree $2m$, namely $H^0(F,\Z_\ell(m))$, vanishes, since the $\ell$-adic cyclotomic character of $F$ has infinite image.
Under geometric Morita equivalence, $N_{D/F}$ acts on these groups as multiplication by $\deg D$, and $\theta_E$ acts as $\cor_{E/F}$.

\emph{Case $\ell\ne p$.}
Let $E/F$ be the unramified extension of degree $\deg D$, which splits $D$.
We first show that
\[
 \cor_{E/F}\colon\pi_j\bigl(K^{\et}(D_E)^\wedge_\ell\bigr)\longrightarrow\pi_j\bigl(K^{\et}(D)^\wedge_\ell\bigr)
\]
is surjective for $j\ge2$.
Let $I_F\subset G_F$ be the inertia subgroup, so that $G_F/I_F\simeq G_k$ and $\cd(G_k)=1$.
Since $\ell\ne p$, we have $H^1(I_F,\Z_\ell(m))\simeq\Z_\ell(m-1)$ as $G_k$-modules.
For $m\ge2$ we have $H^0(k,H^1(I_F,\Z_\ell(m)))=H^0(k,\Z_\ell(m-1))=0$, and hence
\[
 H^1(F,\Z_\ell(m))\simeq H^1(k,\Z_\ell(m))\quad(m\ge2),\qquad
 H^2(F,\Z_\ell(m+1))\simeq H^1(k,\Z_\ell(m))\quad(m\ge1).
\]
On $H^1(k,\Z_\ell(m))$, corestriction along the unramified extension $E/F$ is the quotient map onto Frobenius coinvariants, so it is surjective.
To deduce surjectivity of $\pi_j((c_D)^\wedge_\ell)$ for $j\ge2$, we lift a class along the isomorphism $\pi_j((c_E)^\wedge_\ell)$, where $K(D_E)\simeq K(E)$ by Morita equivalence, and use
\[
 c_D\circ\theta_E=\cor_{E/F}\circ c_E.
\]

For injectivity, suppose first that $p\nmid\deg D$.
Suslin and Yufryakov compute $K_*(D)$ by reduction of the maximal order and localization \cite[Lemma 2 and Theorem 3]{SuslinYufryakov}.
It follows that, for $j\ge2$,
\[
 \#\varprojlim_aK_j(D)/\ell^a=\#\pi_j\bigl(K^{\et}(D)^\wedge_\ell\bigr)<\infty.
\]
Their decomposition and the description of $K_1(D)$ show that $K_{j-1}(D)[\ell^\infty]$ is finite, so its Tate module $\varprojlim_aK_{j-1}(D)[\ell^a]$ in the Bockstein sequence vanishes.
Since the finite-coefficient groups are finite, the Moore towers satisfy the Mittag--Leffler condition and the Milnor $\varprojlim^1$ terms vanish; the surjection of the preceding paragraph is therefore a surjection between finite groups of the same order, hence an isomorphism.
For general $D$, let $F'/F$ be the unramified extension of degree $p^{v_p(\deg D)}$, so that the index of $D_{F'}$ is prime to $p$.
Since $\cor_{F'/F}\circ\res_{F'/F}=[F':F]\in\Z_\ell^\times$, restriction is injective on $\pi_j(K(D)^\wedge_\ell)$, so the case over $F'$ shows that the kernel of the map induced by $(c_D)^\wedge_\ell$ vanishes over $F$.

\emph{Case $\operatorname{char}F=0$ and $\ell=p$.}
In this case we use the trace comparison with closed support of Hesselholt--Larsen--Lindenstrauss \cite[Theorem B(1) and Theorem 2.5]{HesselholtLarsenLindenstrauss2020}, which gives
\[
 \tau_{\ge2}K(D)^\wedge_p\simeq\tau_{\ge2}K(F)^\wedge_p,
\]
also in residue characteristic two by the paragraph following \cite[Theorem B]{HesselholtLarsenLindenstrauss2020}.
Since $L_{K(1),p}HM\simeq0$ for every abelian group $M$, a connective spectrum $T$ satisfies $L_{K(1),p}\tau_{\ge2}T\simeq L_{K(1),p}T$.
Hence, for connective $T$, the localization unit induces an isomorphism on $\pi_j$ for $j\ge2$ if and only if it does so for $\tau_{\ge2}T$.
By the Quillen--Lichtenbaum comparison for $F$ \cite[Theorems 6.13(1) and 7.13--7.14]{ClausenMathew}, the map $K(F)^\wedge_p\to L_{K(1),p}K(F)$ induces an isomorphism on $\pi_j$ for $j\ge2$.
The result for $D$ follows from naturality of the localization unit and the displayed equivalence.
By the factorization of the localization unit above, which holds also for $\ell=p$ in characteristic zero \cite[Example 6.3]{ClausenMathew}, $(c_D)^\wedge_p$ induces an isomorphism on $\pi_j$ for $j\ge2$.

\emph{Case $\operatorname{char}F=\ell=p$.}
Then $[F:F^p]=p$, and Proposition~\ref{prop:app-frobenius}(2) gives
\[
 \pi_j\bigl(K(D)/p^a\bigr)=0\qquad(j\ge2).
\]
Since $[F^{\mathrm{sep}}:(F^{\mathrm{sep}})^p]=p$ as well, the same vanishing holds for $K(D_{F^{\mathrm{sep}}})/p^a$ after split Morita equivalence.
This upper bound is preserved by global sections of the \'etale hypersheaf and by the inverse limit.
\end{proof}


\begin{lemma}\label{lem:paper-local-galois-vanishing}
For $t\ge2$ and $s\ge1$, we have
\[
 H^s\bigl(F,K_t(F^{\mathrm{sep}})\bigr)=0
\]
unless $s=1$ and $t$ is odd.
\end{lemma}
\begin{proof}
For $t\ge2$ the groups $K_t(F^{\mathrm{sep}})$ are divisible, and for $b\ge1$
\[
 K_{2b}(F^{\mathrm{sep}})\otimes\Q\simeq K_{2b}(F^{\mathrm{sep}}),\qquad
 K_{2b+1}(F^{\mathrm{sep}})_{\mathrm{tors}}\simeq\bigoplus_{\ell\ne\operatorname{char}F}\Q_\ell/\Z_\ell(b+1),
\]
where in mixed characteristic the sum includes $\ell=p$, and $K_{2b+1}(F^{\mathrm{sep}})/K_{2b+1}(F^{\mathrm{sep}})_{\mathrm{tors}}$ is a rational vector space.
For an algebraic closure of $F$ this is Suslin's theorem \cite{SuslinClosed}, and in characteristic zero $F^{\mathrm{sep}}$ is algebraically closed.
In characteristic $p$, restriction from $F^{\mathrm{sep}}$ to its algebraic closure is an isomorphism after inverting $p$, since for every finite purely inseparable extension $F'/F^{\mathrm{sep}}$
\[
 N_{F'/F^{\mathrm{sep}}}\res_{F'/F^{\mathrm{sep}}}=\res_{F'/F^{\mathrm{sep}}}N_{F'/F^{\mathrm{sep}}}=[F':F^{\mathrm{sep}}].
\]
For the $p$-part, $F^{\mathrm{sep}}$ and its finite purely inseparable extensions $F'$ satisfy $[F':F'^p]=p$, so $K_t^M(F')/p=0$ for $t\ge2$ by the logarithmic differentials, and the calculation of Geisser and Levine together with the Bockstein sequences in degrees $t$ and $t+1$ shows that $K_t(F^{\mathrm{sep}})$ is uniquely $p$-divisible for $t\ge2$ \cite{GeisserLevine}.

Positive-degree cohomology of discrete rational coefficients vanishes by averaging, which proves the assertion for $t$ even.
For $t=2b+1$ odd and $s\ge2$, the long exact sequence for the torsion subgroup and the rational quotient reduces the assertion to the torsion subgroup.
There $H^s$ vanishes for $s\ge3$ since $\cd_\ell(F)\le2$, and for $s=2$ local Tate duality \cite{TateDuality} and $H^0(F,\Z_\ell(-b))=0$ give
\[
 H^2\bigl(F,\Q_\ell/\Z_\ell(b+1)\bigr)=0.
\]
\end{proof}


For a central simple algebra $B$ over a field $F$, the \emph{geometric rank} of a class in $\pi_0K^{\et}(B)$ is its image in $\pi_0K^{\et}(B_{F^{\mathrm{sep}}})\simeq\pi_0K(F^{\mathrm{sep}})=\Z$ under Morita equivalence; similarly, for $\ell\ne\operatorname{char}F$, the geometric rank of a class in $\pi_0L_{K(1),\ell}K(B)$ is its image in $\pi_0L_{K(1),\ell}K(B_{F^{\mathrm{sep}}})\simeq\pi_0L_{K(1),\ell}K(F^{\mathrm{sep}})=\Z_\ell$.
Classes of geometric rank $1$ are the basic input for the comparisons of Section~\ref{sec:paper-local-spectrum} and Section~\ref{sec:global-spectrum}.

\begin{lemma}\label{lem:paper-local-low-descent}
The group $\pi_0K^{\et}(D)$ is torsion-free, and under the identification $\pi_1K^{\et}(D)\simeq F^\times$ given by the edge map, the map $\pi_1(c_D)$ is the reduced norm isomorphism
\[
 \Nrd_1\colon K_1(D)\xrightarrow{\ \sim\ }F^\times.
\]
Moreover $(c_D)^\wedge_\ell$ induces an isomorphism on $\pi_1$ for every prime $\ell$.
If $\ell\ne\operatorname{char}F$, the localization unit $K(D)^\wedge_\ell\to L_{K(1),\ell}K(D)$ also induces an isomorphism on $\pi_1$.
\end{lemma}
\begin{proof}
We use the integral descent spectral sequence
\[
 E_2^{s,t}=H^s\bigl(F,K_t(F^{\mathrm{sep}})\bigr)\Longrightarrow\pi_{t-s}K^{\et}(D),
\]
whose coefficients are identified with those of $F$ by Morita equivalence.
For any discrete module $W$, we first separate its torsion subgroup and then rationalize the torsion-free quotient; since $\cd_\ell(F)\le2$ for torsion coefficients, we get $H^s(F,W)=0$ for $s>3$, so the spectral sequence converges.
By Lemma~\ref{lem:paper-local-galois-vanishing}, in the rows $t\ge2$ the only possible nonzero terms with $s\ge1$ are $H^1(F,K_t(F^{\mathrm{sep}}))$ with $t$ odd.
In total degree one, $H^1(F,K_2(F^{\mathrm{sep}}))$ and $H^2(F,K_3(F^{\mathrm{sep}}))$ therefore vanish, as do the targets $E_r^{r,r}$ of the differentials leaving $E_r^{0,1}$.
Thus the edge map gives $\pi_1K^{\et}(D)\simeq F^\times$.
The composite of $\pi_1(c_D)$ with this edge map extends scalars to $D_{F^{\mathrm{sep}}}\simeq M_{\deg D}(F^{\mathrm{sep}})$ and then takes the determinant, so it is $\Nrd_1$.
The map $\Nrd_1\colon K_1(D)\to F^\times$ is surjective: $D$ has ramification index $\deg D$ over $F$ and contains an unramified maximal subfield, so $\Nrd_1$ takes values of every valuation, and the norm from the unramified maximal subfield is surjective on units \cite[V, \S2, Proposition 3(b) and its Corollary; XII, \S2, Propositions 1--2 and Exercises 2--3]{SerreLocal}.
It is injective since $\SK_1(D)=0$ \cite[Corollary 4.10(iii)]{HazratWadsworth}; the hypotheses there hold because $D$ is tame and semiramified, with residue field cyclic of degree $\deg D$ over $k$, and the Tate group $\widehat H^{-1}$ of the cyclic residue Galois group appearing there vanishes by Hilbert 90.
In total degree zero,
\[
 H^1\bigl(F,(F^{\mathrm{sep}})^\times\bigr)=0,\qquad H^2\bigl(F,K_2(F^{\mathrm{sep}})\bigr)=0,
\]
by Hilbert 90 and by Lemma~\ref{lem:paper-local-galois-vanishing}, respectively.
Hence the geometric rank map $\pi_0K^{\et}(D)\to H^0(F,\Z)=\Z$ is injective, and $\pi_0K^{\et}(D)$ is torsion-free.

Let $T$ be $K(D)$ or $K^{\et}(D)$.
Since $\pi_0T$ is torsion-free, the Bockstein sequence gives $\pi_1(T/\ell^a)\simeq\pi_1(T)/\ell^a$.
The Milnor sequence for the derived completion is
\[
 0\to\varprojlim\nolimits^1_a\pi_2(T/\ell^a)
 \to\pi_1(T^\wedge_\ell)\to\varprojlim_a\pi_1(T/\ell^a)\to0.
\]
The transition maps of $\{\pi_2(T)/\ell^a\}_a$ are surjective, so the derived inverse limit of the degree-two Bockstein sequence gives
\[
 \varprojlim\nolimits^1_a\pi_2(T/\ell^a)\simeq\varprojlim\nolimits^1_a\pi_1(T)[\ell^a].
\]
The isomorphism $\pi_1(c_D)$ induces isomorphisms on both outer terms, which proves the claim in degree one.

Let $\ell\ne\operatorname{char}F$.
Then $\pi_1L_{K(1),\ell}K(D)\simeq H^1(F,\Z_\ell(1))$, and Kummer theory and Hilbert 90 give $H^1(F,\mu_{\ell^a})\simeq F^\times/(F^\times)^{\ell^a}$.
Since $F^\times[\ell^\infty]$ is finite, the inverse system $\{F^\times[\ell^a]\}_a$ of $\ell$-power roots of unity is Mittag--Leffler, and $\varprojlim^1_aF^\times[\ell^a]=0$.
Thus
\[
 H^1(F,\Z_\ell(1))\simeq\varprojlim_a F^\times/(F^\times)^{\ell^a}.
\]
Together with $\pi_1(c_D)=\Nrd_1$, these isomorphisms show that $K(D)^\wedge_\ell\to L_{K(1),\ell}K(D)$ induces an isomorphism on $\pi_1$ as well.
\end{proof}


\begin{lemma}\label{lem:paper-local-compare}
The \'etale comparison map $c_D\colon K(D)\to K^{\et}(D)$ induces an isomorphism on $\pi_j$ for every $j\ge1$.
For $\ell\ne\operatorname{char}F$, the localization unit $K(D)^\wedge_\ell\to L_{K(1),\ell}K(D)$ also induces an isomorphism on $\pi_j$ for every $j\ge1$.
\end{lemma}
\begin{proof}
Set $Y=\cofib(c_D)$.
Let $E/F$ be a finite separable splitting field and $Y_E=\cofib(c_{D_E})$.
By Morita equivalence and the comparison for fields above, $\pi_iY_E=0$ for $i\ge1$.
Since $\cor_{E/F}\circ\res_{E/F}=[E:F]$ on $\pi_iY$, we obtain
\[
 [E:F]\cdot\pi_iY=0\qquad(i\ge1).
\]
By Lemmas~\ref{lem:paper-local-complete-high} and~\ref{lem:paper-local-low-descent}, $\pi_i(Y^\wedge_\ell)=0$ for $i\ge2$ and $\pi_1(c_D)=\Nrd_1$ is an isomorphism.
The Bockstein sequence for $Y/\ell\simeq Y^\wedge_\ell/\ell$ then gives $\pi_i(Y)[\ell]=0$ for $i\ge2$, and the bounded exponent implies $\pi_iY=0$ for $i\ge2$.
For $\ell\ne\operatorname{char}F$, the same two lemmas give the height-one comparison on $\pi_j$ for $j\ge2$ and on $\pi_1$, respectively.
\end{proof}

To define the even-degree isomorphisms of Theorem~\ref{thm:paper-local-norm}, we describe splitting transfers from Galois extensions in terms of coinvariants.

\begin{lemma}\label{lem:paper-local-even-coinvariants}
Let $L/F$ be a finite Galois extension with group $G$, and let $m\ge1$.
Then the norm $N_{D_L/D}$ induces an isomorphism
\[
 K_{2m}(D_L)_G\xrightarrow{\ \sim\ }K_{2m}(D).
\]
Moreover $K_{2m}(D)/K_{2m}(D)_{\mathrm{tors}}$ is a rational vector space, and $N_{D_L/D}$ maps $(K_{2m}(D_L)_{\mathrm{tors}})_G$ onto $K_{2m}(D)_{\mathrm{tors}}$.
All these assertions hold with $D$ replaced by $F$.
\end{lemma}
\begin{proof}
Set
\[
 M=K_{2m+1}(F^{\mathrm{sep}}),\qquad P=K_{2m}(F^{\mathrm{sep}}).
\]
By integral descent, Lemma~\ref{lem:paper-local-compare} over $L$, and Lemma~\ref{lem:paper-local-galois-vanishing}, we have a natural exact sequence
\[
 0\longrightarrow H^1(L,M)\longrightarrow K_{2m}(D_L)
 \longrightarrow H^0(L,P)\longrightarrow0.
\]
The right-hand term is a rational vector space and the left-hand term is torsion, so $H^1(L,M)$ is the torsion subgroup of $K_{2m}(D_L)$.
For $L=F$ this shows that $K_{2m}(D)/K_{2m}(D)_{\mathrm{tors}}$ is a rational vector space.
Restriction and transfer are compatible with these exact sequences and act on the outer terms by Galois restriction and corestriction, since at geometric points they are the diagonal map and the sum over the factors of the separable base change, respectively.
Under local Tate duality, the dual of the corestriction
\[
 H^1(L,\Q_\ell/\Z_\ell(m+1))_G\longrightarrow H^1(F,\Q_\ell/\Z_\ell(m+1))
\]
is restriction \cite{TateDuality}.
Since the dual coefficients are $\Z_\ell(-m)$ and $H^0(L,\Z_\ell(-m))=0$, inflation--restriction identifies this restriction map with an isomorphism onto the invariants.
For $M/M_{\mathrm{tors}}$ and $P$ we use averaging over rational coefficients: on rational vector spaces with $G$-action, the norm induces an isomorphism from coinvariants to invariants.
We have
\[
 H^1(L,M)\simeq\coker\bigl(H^0(L,M/M_{\mathrm{tors}})\to H^1(L,M_{\mathrm{tors}})\bigr),\qquad H_1(G,H^0(L,P))=0,
\]
so corestriction induces isomorphisms on the coinvariants of the two outer terms, and hence on $K_{2m}(D_L)_G$.
Since the left-hand terms are the torsion subgroups, the isomorphism on their coinvariants, $(K_{2m}(D_L)_{\mathrm{tors}})_G\xrightarrow[\ \sim\ ]{N_{D_L/D}}K_{2m}(D)_{\mathrm{tors}}$, is the second assertion.
The same calculation applies to $F$.
\end{proof}

\begin{proof}[Proof of Theorem~\ref{thm:paper-local-norm}]
We first treat odd degrees $j=2m-1\ge3$.
By Lemma~\ref{lem:paper-local-galois-vanishing}, the only contribution to total degree $j$ in descent is $H^0(F,K_j(F^{\mathrm{sep}}))$, and the differentials leaving it land in vanishing groups.
By Lemma~\ref{lem:paper-local-compare} applied to $D$ and to $F$, the \'etale comparison maps give $K_j(D)\simeq\pi_jK^{\et}(D)$ and $K_j(F)\simeq\pi_jK^{\et}(F)$.
The edge maps
\[
 e_D\colon\pi_jK^{\et}(D)\longrightarrow H^0(F,K_j(F^{\mathrm{sep}})),\qquad
 e_F\colon\pi_jK^{\et}(F)\longrightarrow H^0(F,K_j(F^{\mathrm{sep}}))
\]
identify the groups and are compatible with restriction and transfer.
We define $\nu_{j,D}$ as the composite
\[
 K_j(D)\xrightarrow{\ \pi_j(c_D)\ }\pi_jK^{\et}(D)\xrightarrow{\ e_D\ }H^0(F,K_j(F^{\mathrm{sep}}))\xrightarrow{\ e_F^{-1}\ }\pi_jK^{\et}(F)\xrightarrow{\ \pi_j(c_F)^{-1}\ }K_j(F).
\]
For the normalization $\nu_{j,D}\theta_E=N_{E/F}$, the splitting field $E$ need not be separable.
Let $E_s$ be its maximal separable subextension.
The algebra $E\otimes_FF^{\mathrm{sep}}$ is a product, indexed by the embeddings of $E_s$, of purely inseparable field extensions of $F^{\mathrm{sep}}$, over each of which $D$ splits.
Under the common geometric Morita identification, the component of $\theta_E$ at each factor is the corresponding field norm.
Flat base change and injectivity of the edge map therefore give $\nu_{j,D}\theta_E=N_{E/F}$.
Under the common geometric Morita identification for $e_D$ and $e_F$, restriction of scalars at a geometric point is multiplication by $\deg D$, so injectivity of $e_F\circ\pi_j(c_F)$ gives $\deg D\cdot\nu_{j,D}=\pi_jN_{D/F}$ in odd degrees.

Next let $j=2m\ge2$ be even, and let $L/F$ be a finite Galois splitting field with group $G$.
By Lemma~\ref{lem:paper-local-even-coinvariants}, the maps $\theta_L$ and $N_{L/F}$ induce isomorphisms
\[
 K_{2m}(D)\xleftarrow[\ \sim\ ]{\ \theta_L\ }K_{2m}(L)_G\xrightarrow[\ \sim\ ]{\ N_{L/F}\ }K_{2m}(F),
\]
and we define $\nu_{2m,D}=N_{L/F}\circ\theta_L^{-1}$, where $\theta_L$ and $N_{L/F}$ denote the induced maps on coinvariants.
The induced identification $K_{2m}(L)\simeq K_{2m}(D_L)$ is $G$-equivariant, since conjugation takes the defining simple bimodule to an isomorphic one.
This definition does not depend on $L$: two Galois splitting fields can be compared over their compositum, using the surjectivity of the field norm from the compositum and the transitivity of transfer.
For any finite splitting field $E/F$, the extension $LE/E$ is Galois, so its even-degree field norm is surjective by Lemma~\ref{lem:paper-local-even-coinvariants} applied over $E$, and transitivity of transfer gives $\nu_{2m,D}\theta_E=N_{E/F}$.
For a finite separable extension $F'/F$, choose a Galois splitting field $L/F$ containing $F'$; then $\nu_{2m,D}$ is compatible with restriction and transfer along $F'/F$.
Indeed, the Mackey formula for restriction of tensor products gives the restriction identity on the image of $\theta_L$, and transitivity of transfer gives the transfer identity on the image of the splitting transfer $\theta_{L/F'}\colon K(L)\to K(D_{F'})$ for $D_{F'}$.
The identity $N_{D/F}\theta_L=\deg D\cdot N_{L/F}$ gives $\deg D\cdot\nu_{2m,D}=\pi_{2m}N_{D/F}$.
In degree two, the transfer formula $\Nrd_2\theta_L=N_{L/F}$ \cite[Proposition 6.6.1(1)]{KahnLevine2010} and the surjectivity of $\theta_L$ give $\nu_{2,D}=\Nrd_2$.

It remains to treat degree one.
We set $\nu_{1,D}=\Nrd_1$, which is an isomorphism by Lemma~\ref{lem:paper-local-low-descent}.
Compatibility with restriction follows from base change for reduced norms.
The transfer formula for reduced norms \cite[Proposition 6.6.1(1)]{KahnLevine2010} gives $\Nrd_1\theta_E=N_{E/F}$ and $\Nrd_1\circ N_{D_{F'}/D}=N_{F'/F}\circ\Nrd_1$.
For $x\in D^\times$, the determinant of multiplication by $x$ on the $F$-vector space $D$ is $\Nrd_1(x)^{\deg D}$, so
\[
 \pi_1N_{D/F}=\deg D\cdot\Nrd_1\quad\text{on }K_1(D).
\]
\end{proof}
